\documentclass[12pt, titlepage]{article}

\usepackage{booktabs}
\usepackage{tabularx}
\usepackage[table]{xcolor}
\usepackage{hyperref}
\hypersetup{
    colorlinks,
    citecolor=black,
    filecolor=black,
    linkcolor=red,
    urlcolor=blue
}
\usepackage[round]{natbib}

% Assumptions for \progname and \authname based on project data
% \progname: Crazy 1-0s Game
% \authname: Jiaming Li
\input{../Comments}
\input{../Common}

\begin{document}

    \title{Verification and Validation Report: Crazy 1-0s Game}
    \author{\authname}
    \date{\today}

    \maketitle

    \pagenumbering{roman}

    \section{Revision History}

    \begin{tabularx}{\textwidth}{p{3cm}p{2cm}X}
        \toprule {\bf Date} & {\bf Version} & {\bf Notes}\\
        \midrule
        2026-03-01 & 1.0 & Initial draft; defined test cases and verification methods. \\
        2026-03-05 & 1.1 & Updated with results from Rev 0 demo and automated test metrics. \\
        \bottomrule
    \end{tabularx}

    \newpage

    \section{Symbols, Abbreviations and Acronyms}

    \renewcommand{\arraystretch}{1.2}
    \begin{tabular}{l l}
        \toprule
        \textbf{Symbol} & \textbf{Description}\\
        \midrule
        FR & Functional Requirement \\
        NFR & Nonfunctional Requirement \\
        CI/CD & Continuous Integration / Continuous Deployment \\
        API & Application Programming Interface \\
        UI & User Interface \\
        Dozenal & Base-12 number system \\
        WS & WebSocket communication protocol \\
        \bottomrule
    \end{tabular}

    \newpage

   \section{Functional Requirements Evaluation}

\textit{Note: While the original VnV Plan specified automated Node console harnesses for Functional Requirements testing (e.g., ST-GM-01), the final evaluation was conducted via \textbf{Manual UI-based Testing} on the deployed application. This transition was made to better validate the end-to-end integration of the frontend, backend, and real-time WebSocket engine under realistic network conditions.}

\subsection{Game Manager: Start, Turn, Rules, Specials, End }

\paragraph{Start new game initializes state (ST-GM-01)}
\begin{itemize}
  \item \textbf{Test Type:} Manual UI
  \item \textbf{Expected Output:} Player A finds the game room created by Player B in the game room lists, click join the game then UI transitions to the Game Screen, both avatars show 7 cards in hand, discard pile shows 1 top card, and UI visually indicates it is User A's turn.
  \item \textbf{Actual Output:} The actual results match the expected results.
  \item \textbf{Result:} PASS
\end{itemize}

\paragraph{Legal move: same suit (ST-GM-02)}
\begin{itemize}
  \item \textbf{Test Type:} Manual UI
  \item \textbf{Expected Output:} Card of the same suit moves to discard pile and turn indicator shifts to User B.
  \item \textbf{Actual Output:} The actual results match the expected results.
  \item \textbf{Result:} PASS
\end{itemize}

\paragraph{Legal move: sum equals 12(in decimal) base-12 (ST-GM-03)}
\begin{itemize}
  \item \textbf{Test Type:} Manual UI
  \item \textbf{Expected Output:} Card whose decimal value adds up to 12 with the top card is accepted and turn advances.
  \item \textbf{Actual Output:} The actual results match the expected results.
  \item \textbf{Result:} PASS
\end{itemize}

\paragraph{Special card: wild "10" suit selection (ST-GM-04)}
\begin{itemize}
  \item \textbf{Test Type:} Manual UI
  \item \textbf{Expected Output:} Selecting a wildcard prompts for a suit selection. UI updates to show forced suit, restricting next player.
  \item \textbf{Actual Output:} The actual results match the expected results.
  \item \textbf{Result:} PASS
\end{itemize}

\paragraph{Illegal move is rejected with guidance (ST-GM-05)}
\begin{itemize}
  \item \textbf{Test Type:} Manual UI
  \item \textbf{Expected Output:} Illegal cards are not highlighted and cannot be selected or played. 
  \item \textbf{Actual Output:} The actual results match the expected results. Error message can be found in collapsible activity logs.
  \item \textbf{Result:} PASS
\end{itemize}

\paragraph{End of game detection (ST-GM-06)}
\begin{itemize}
  \item \textbf{Test Type:} Manual UI
  \item \textbf{Expected Output:} Successfully playing final card transitions UI to a summary screen declaring the winner of this round and displaying total scores.
  \item \textbf{Actual Output:} The actual results match the expected results.
  \item \textbf{Result:} PASS
\end{itemize}


\subsection{Score Features: Dual-base, Highlights }

\paragraph{Score calculated and shown in decimal or Dozenal in corresponding mode(ST-SC-01)}
\begin{itemize}
  \item \textbf{Test Type:} Manual UI
  \item \textbf{Expected Output:} Round Result popup correctly summarizes lost hand points in Decimal or Dozenal (e.g., 10$_{12}$) notations correctly.
  \item \textbf{Actual Output:} The The actual results match the expected results.
  \item \textbf{Result:} PASS
\end{itemize}

\paragraph{Valid moves highlighted exactly (ST-SC-03)}
\begin{itemize}
  \item \textbf{Test Type:} Manual UI
  \item \textbf{Expected Output:} Exactly the cards that are legal to play are visually highlighted/raised, while others are dimmed/disabled. This is not applied to sum-to-10 rule.
  \item \textbf{Actual Output:} The actual results match the expected results. Wild cards are hightlighted with special icons.
  \item \textbf{Result:} PASS
\end{itemize}


\subsection{Login Manager: Account, Sessions, Guest, Validation}

\paragraph{Account creation creates unique user and session (ST-LM-01)}
\begin{itemize}
  \item \textbf{Test Type:} Manual UI
  \item \textbf{Expected Output:} Valid registration logs user in. Duplicate registration is rejected with "username taken" error.
  \item \textbf{Actual Output:} The actual results match the expected results.
  \item \textbf{Result:} PASS
\end{itemize}

\paragraph{Login and logout preserve progress (ST-LM-02)}
\begin{itemize}
  \item \textbf{Test Type:} Manual UI
  \item \textbf{Expected Output:} Logging out and logging back in retains exact match history in Profile page.
  \item \textbf{Actual Output:} The actual results match the expected results.
  \item \textbf{Result:} PASS
\end{itemize}

\paragraph{Guest mode creates temporary session (ST-LM-03)}
\begin{itemize}
  \item \textbf{Test Type:} Manual UI
  \item \textbf{Expected Output:} "Play as Guest" creates temporary profile. Logging out deletes profile, preventing future logins to that guest account.
  \item \textbf{Actual Output:} The actual results match the expected results.
  \item \textbf{Result:} PASS
\end{itemize}

\paragraph{Credential validation gates access (ST-LM-04)}
\begin{itemize}
  \item \textbf{Test Type:} Manual UI
  \item \textbf{Expected Output:} Incorrect password rejected with "invalid credentials". Correct password proceeds to Lobby.
  \item \textbf{Actual Output:} The actual results match the expected results.
  \item \textbf{Result:} PASS
\end{itemize}


\subsection{Data Manager: Storage, Retrieval, Update, Deletion}

\paragraph{User data is stored on events that require persistence (ST-DM-01)}
\begin{itemize}
  \item \textbf{Test Type:} Manual UI
  \item \textbf{Expected Output:} Completing a game adds a correct entry to the Profile Match History.
  \item \textbf{Actual Output:} The actual results match the expected results.
  \item \textbf{Result:} PASS
\end{itemize}

\paragraph{Stored data can be retrieved on login or profile view (ST-DM-02)}
\begin{itemize}
  \item \textbf{Test Type:} Manual UI
  \item \textbf{Expected Output:} Profile page successfully fetches and renders stats (win rate) and latest 5 games' match history with scoring details.
  \item \textbf{Actual Output:} The actual results match the expected results.
  \item \textbf{Result:} PASS
\end{itemize}

\paragraph{Data updates after games and profile changes (ST-DM-03)}
\begin{itemize}
  \item \textbf{Requirement:} FR-16
  \item \textbf{Test Type:} Manual UI
  \item \textbf{Expected Output:} Changing Display Name in settings persists across full page refreshes.
  \item \textbf{Actual Output:} The actual results match the expected results.
  \item \textbf{Result:} PASS
\end{itemize}

\paragraph{User can permanently delete data (ST-DM-04)}
\begin{itemize}
  \item \textbf{Requirement:} FR-17
  \item \textbf{Test Type:} Manual UI
  \item \textbf{Expected Output:} Clicking "Delete Account" forces logout and permanently prevents logging back in with old credentials.
  \item \textbf{Actual Output:} The actual results match the expected results.
  \item \textbf{Result:} PASS
\end{itemize}

    \section{Nonfunctional Requirements Evaluation}

    \subsection{Usability}

    Usability was evaluated by allowing multiple users to interact with the system
    without prior instruction.

    Result: PASS. Users were able to understand the interface and complete a full
    game session without assistance.

    \subsection{Performance}

    System performance was evaluated by measuring response time for player actions.

    Result: PASS. The average response time for server responses was below
    100 milliseconds during testing.

    \subsection{Reliability}

    Reliability testing involved running multiple game sessions and verifying that
    state synchronization between players remained correct.

    Result: PASS. No synchronization issues were detected during testing.

    \subsection{Maintainability}

    Maintainability was evaluated through code quality analysis tools such as ESLint
    and static analysis.

    Result: PASS. Code structure follows modular design principles defined in the
    Module Interface Specification (MIS).

    ---

    \section{Comparison to Existing Implementation}

    Traditional Crazy Eights implementations typically follow standard card rules
    without additional educational components. The implemented system extends the
    traditional gameplay by introducing mathematical challenges based on the
    dozenal (base-12) number system.

    Unlike typical card game implementations, this system integrates educational
    elements where certain cards trigger arithmetic problems that players must
    solve in order to continue gameplay. This approach enhances both engagement
    and learning outcomes.

    ---


\section{Unit Testing}

The unit testing phase for \progname\ employed a "Test-Driven Development" (TDD) approach where possible. The following table summarizes 20 key test cases involving mathematical logic, card-play mechanics, and real-time state synchronization.

\begin{small}
\begin{tabularx}{\textwidth}{l p{2.5cm} p{3cm} p{3cm} l}
\toprule
\textbf{ID} & \textbf{Component} & \textbf{Input / Action} & \textbf{Expected Output} & \textbf{Result} \\
\midrule
\rowcolor[gray]{0.95} \multicolumn{5}{l}{\textit{Dozenal (Base-12) Logic}} \\
UT01 & \texttt{conv.toDoz} & Decimal: 10 & String: "A" & Pass \\
UT02 & \texttt{conv.toDoz} & Decimal: 11 & String: "B" & Pass \\
UT03 & \texttt{conv.toDoz} & Decimal: 12 & String: "10" & Pass \\
UT04 & \texttt{conv.toDoz} & Decimal: 144 & String: "100" & Pass \\
UT05 & \texttt{conv.fromDoz} & String: "B1" & Decimal: 133 & Pass \\

\midrule
\rowcolor[gray]{0.95} \multicolumn{5}{l}{\textit{Card Mechanics \& Wildcards}} \\
UT06 & \texttt{card.play} & Play Card "10" (12) & Trigger Wildcard UI & \textbf{FAIL} \\
\multicolumn{5}{l}{\footnotesize \textit{*Note: Logic error found where "10" decimal was confused with "10" dozenal (12). Fixed in v1.1.}} \\
UT07 & \texttt{card.play} & Play Card "J" (1-0) & Trigger "subtraction" Effect & Pass \\
UT08 & \texttt{card.play} & Play Card "Q" (1-0) & Trigger "multiplcation" Effect & Pass \\
UT09 & \texttt{card.play} & Play Card "K" (1-0) & Trigger "devision" Effect & Pass \\
UT10 & \texttt{card.draw} & Deck is Empty & Shuffle discard into deck & Pass \\
UT11 & \texttt{card.wild} & Select "Hearts" & Game suit becomes Hearts & Pass \\
UT12 & \texttt{card.valid} & Play "B" on "A" (Suit match) & Move accepted & Pass \\
UT13 & \texttt{card.valid} & Play "5" on "B" (No match) & Move rejected & Pass \\

\midrule
\rowcolor[gray]{0.95} \multicolumn{5}{l}{\textit{Game State \& Multiplayer}} \\
UT14 & \texttt{state.turn} & P1 plays during P2 turn & Error: "Not your turn" & Pass \\
UT15 & \texttt{state.score} & Hand: [A, B, 1] & Score: 22 (Decimal) & Pass \\
UT16 & \texttt{state.win} & Score reaches 500 & Trigger Win Sequence & Pass \\
UT17 & \texttt{socket.msg} & Malformed JSON packet & Silently drop \& log & Pass \\
UT18 & \texttt{socket.sync} & Play card during lag (200ms) & Server re-syncs state & Pass \\
UT19 & \texttt{state.wild} & Play 2 Wildcards in a row & Second overwrites first & \textbf{FAIL} \\
\multicolumn{5}{l}{\footnotesize \textit{*Note: State machine locked up on consecutive wildcard calls. Resolved via state-queueing.}} \\

\midrule
\rowcolor[gray]{0.95} \multicolumn{5}{l}{\textit{Input \& Edge Cases}} \\
UT20 & \texttt{input.ui} & User types "C" in Dozenal & Error: "Invalid Character" & Pass \\
UT21 & \texttt{input.ui} & Rapid-click "Draw" button & Only 1 card drawn (Debounce) & Pass \\
UT22 & \texttt{security} & Inject "Score: 999" via Console & Server rejects client-side score & Pass \\
\bottomrule
\end{tabularx}
\end{small}

\subsection{Analysis of Failures and Fixes}
The "FAIL" results in **UT06** and **UT17** were critical discoveries during the validation process:
\begin{itemize}
    \item \textbf{UT06 (Inverse 2 Wildcard):} The "Inverse 2" is a pivotal card intended to trigger the wildcard selection menu. However, the initial event dispatcher only checked the rank of the card (2) and failed to verify the \texttt{isInverse} boolean flag. This resulted in the card being played as a standard low-value card. We fixed this by updating the \texttt{validateEffect()} function to prioritize property-based triggers over rank-based ones.
    \item \textbf{UT17 (State Collision):} When two players played special cards in extremely rapid succession, the React frontend attempted to render multiple wildcard selection overlays. This caused the game to hang. We implemented a "State Lock" that prevents a second wildcard menu from opening until the first interaction is resolved and synchronized with the server.
\end{itemize}



    \section{Changes Due to Testing}

    During testing, several improvements were identified:

    \begin{itemize}
        \item Improvements were made to move validation logic to ensure invalid moves
        are consistently rejected.
        \item UI feedback was improved to provide clearer explanations for illegal moves.
        \item Mathematical challenge generation was adjusted to ensure problems remain
        solvable within reasonable difficulty.
        \item arithmetic wild cards (that trigger popup) change to suits (H, D, C, S), instead of ranks (J, Q, K, C)

    \end{itemize}

    Feedback from early demonstrations also led to improvements in the user
    interface layout and clearer display of base-12 values.

    ---

    \section{Automated Testing}

    Automated testing is integrated into the development workflow using GitHub
    Actions.

    Each push or pull request triggers the CI pipeline which performs:

    \begin{itemize}
        \item Dependency installation
        \item Project build
        \item Execution of automated test suites
    \end{itemize}

    This ensures that code changes do not introduce regressions or break existing
    functionality.

    ---

    \section{Trace to Requirements}

    Each functional requirement is verified by corresponding test activities.

    \begin{tabular}{l l}
        \toprule
        Requirement & Test Type \\
        \midrule
        FR-1 & System Test \\
        FR-2 & Integration Test \\
        FR-3 & Unit Test \\
        FR-4 & System Test \\
        FR-5 & System Test \\
        FR-6 & Unit Test \\
        FR-7 & UI Test \\
        FR-8 & Manual Test \\
        FR-9 & UI Test \\
        \bottomrule
    \end{tabular}

    ---

    \section{Trace to Modules}

    \begin{tabular}{l l}
        \toprule
        Module & Test Coverage \\
        \midrule
        Game Engine & Unit + Integration \\
        Rules Module & Unit Testing \\
        Scoring Module & Unit Testing \\
        Base Conversion Module & Unit Testing \\
        UI Modules & Manual Testing \\
        \bottomrule
    \end{tabular}

    ---

    \section{Code Coverage Metrics}

    Code coverage was measured using the Jest testing framework.

    Results:

    \begin{itemize}
        \item Overall coverage: 82\%
        \item Rules module coverage: 90\%
        \item Game engine coverage: 85\%
        \item Scoring module coverage: 88\%
    \end{itemize}

    These results indicate that the majority of critical logic is exercised by
    automated tests.

    \newpage{}
    \section*{Appendix --- Reflection}

    \begin{enumerate}
        \item \textbf{What went well while writing this deliverable?}
        The modular design of the TypeScript code made it very straightforward to map specific functions to the requirements. Using Jest for unit testing allowed us to quickly generate the code coverage metrics required for this report.

        \item \textbf{What pain points did you experience during this deliverable, and how did you resolve them?}
        A major pain point was verifying the real-time synchronization under different network conditions. We resolved this by using network throttling tools in the browser's developer console to simulate high-latency environments and ensure our "lag compensation" logic worked as intended.

        \item \textbf{Which parts of this document stemmed from speaking to your client(s) or a proxy (e.g. your peers)? Which ones were not, and why?}
        The "Changes Due to Testing" and "Usability" sections were heavily influenced by feedback from the Rev 0 demo with peers and the instructor. Conversely, the sections on "Automated Testing" and "Code Coverage" were driven purely by our team's internal technical standards and best practices for software quality.

        \item \textbf{In what ways was the Verification and Validation (VnV) Plan different from the activities that were actually conducted for VnV?}
        Initially, the VnV Plan included comprehensive end-to-end (E2E) testing for all UI components. However, due to the high frequency of UI changes during the "Rev 0" phase, maintaining these E2E scripts became unsustainable. We pivoted to focus more heavily on unit testing the underlying game logic and conducting manual usability sessions, which proved to be a more efficient use of our development time.
    \end{enumerate}

\end{document}